Le chapitre s'ouvre sur le renversement le plus net du livre.

Au lycée, l'intégrale d'une fonction continue positive \emph{est} l'aire sous la courbe :
c'est une définition, et l'existence de cette aire est admise. Ici, l'intégrale est
construite par la théorie de la mesure, et l'aire est la mesure d'une partie du plan. Les deux
points de vue ne se rejoignent donc pas par convention mais par un théorème — le premier du
chapitre — qui établit que la mesure de la partie comprise entre l'axe et la courbe vaut
l'intégrale. C'est l'énoncé que le lycée ne peut qu'admettre, et il est démontré.

Le reste suit sans résistance : primitive définie par une intégrale, unicité à constante près,
théorème fondamental, linéarité, relation de Chasles, positivité, inégalité de la moyenne,
intégration par parties. On notera que la relation de Chasles ne demande aucune condition
d'ordre sur les trois bornes : c'est tout l'intérêt de l'intégrale orientée.

Un énoncé n'est pas traité : le volume d'un solide de révolution. Il faudrait mesurer dans
l'espace la partie engendrée par la rotation d'une courbe, donc disposer du théorème de Fubini
en dimension trois et de l'aire du disque — un appareillage sans rapport avec l'approche du
programme, qui empile des cylindres et passe à la limite sans le justifier.
